\documentclass[acmsmall,nonacm,screen]{acmart}

\usepackage{mathpartir}
\usepackage{booktabs}

\citestyle{acmauthoryear}
\settopmatter{printacmref=false}

\AtEndPreamble{\theoremstyle{acmdefinition}\newtheorem{remark}[theorem]{Remark}}

\newcommand{\kw}[1]{\mathsf{#1}}
\newcommand{\nm}[1]{\mathit{#1}}
\newcommand{\Base}{\kw{Base}}
\newcommand{\Plan}[1]{\kw{Plan}(#1)}
\newcommand{\ret}{\kw{ret}\;}
\newcommand{\model}{\kw{model}}
\newcommand{\letin}[3]{\kw{let}\;#1 = #2\;\kw{in}\;#3}
\newcommand{\ite}[3]{\kw{if}\;#1\;\kw{then}\;#2\;\kw{else}\;#3}
\newcommand{\ask}[1]{\kw{ask}\;#1}
\newcommand{\act}[5]{\kw{act}\;#1\;#2\;#3\;\kw{as}\;(#4)\;\kw{in}\;#5}
\newcommand{\splt}[3]{\kw{split}\;#1\;\kw{as}\;#2 \parallel #3}
\newcommand{\run}[4]{\kw{run}_{#1 \multimap #2}(#3, #4)}
\newcommand{\iter}[4]{\kw{iter}\;#1\;\kw{from}\;#2\;\kw{as}\;#3.\,#4}
\newcommand{\inl}{\kw{inl}\,}
\newcommand{\inr}{\kw{inr}\,}
\newcommand{\tok}[1]{\kw{tok}(#1)}
\newcommand{\sem}[1]{[\![#1]\!]}
\newcommand{\auth}[1]{\mathcal{A}(#1)}
\newcommand{\lang}[1]{\mathcal{L}(#1)}
\newcommand{\sh}{\mathbin{\parallel}}
\newcommand{\wrong}{\kw{wrong}}
\newcommand{\dom}{\mathrm{dom}\,}
\newcommand{\evto}{\Downarrow}
\newcommand{\lstep}[1]{\xrightarrow{\,#1\,}}
\newcommand{\own}{\mathrel{\triangleright}}
\newcommand{\er}[1]{|#1|}
\newcommand{\least}[2]{\lfloor #1 : #2 \rfloor}
\newcommand{\edit}{\nm{edit}}
\newcommand{\test}{\nm{test}}
\newcommand{\push}{\nm{push}}

% A display for programs. It measures its content and warns when the content is wider than the line.
\newsavebox{\progbox}
\newenvironment{prog}
  {\begin{lrbox}{\progbox}$\displaystyle\begin{array}{@{}l@{\qquad}l@{}}}
  {\end{array}$\end{lrbox}%
   \ifdim\wd\progbox>\linewidth
     \PackageWarning{paper}{Overfull program display by \the\dimexpr\wd\progbox-\linewidth\relax}%
   \fi
   \[\usebox{\progbox}\]\ignorespacesafterend}

% written by plans.py; do not edit
\newcommand{\exAccepted}{2,207}
\newcommand{\exActs}{28,985}
\newcommand{\exAsks}{23,196}
\newcommand{\exEnginePlans}{150}
\newcommand{\exEngineWords}{6,000}
\newcommand{\exGrants}{18,561}
\newcommand{\exLent}{699}
\newcommand{\exLucky}{872}
\newcommand{\exNarrowings}{2,849}
\newcommand{\exPairs}{2,000}
\newcommand{\exPairsNo}{1,858}
\newcommand{\exPairsYes}{142}
\newcommand{\exPrograms}{6,000}
\newcommand{\exQuestions}{22,111}
\newcommand{\exRules}{10}
\newcommand{\exRuns}{18,000}
\newcommand{\exSlip}{0.1}
\newcommand{\exSlipBroken}{0.15}
\newcommand{\exSlowest}{7,424}
\newcommand{\exSplits}{3,116}
\newcommand{\exTryAct}{1}
\newcommand{\exTryBranch}{7,424}
\newcommand{\exTryContraction}{102}
\newcommand{\exTryIfFact}{69}
\newcommand{\exTryLoop}{4,450}
\newcommand{\exTryNarrow}{48}
\newcommand{\exTryResidual}{236}
\newcommand{\exTryRun}{12}
\newcommand{\exTryShadow}{535}
\newcommand{\exTrySplit}{33}
\newcommand{\exTyped}{3,000}
\newcommand{\exWritten}{4,347}


\begin{document}

\title{A Plan Is a Language: Typed Approvals for What Comes Next}

\author{Claude-Fable-5.2}
\authornote{An AI model. The proofs are on paper and have not been checked by a machine. A
script implements the calculus and tests the theorems on random programs. At the time of
writing nobody has reviewed the paper.}

\begin{abstract}
An agent that is given a task of any size first proposes a plan, and a person approves it.
Today the plan is prose and the approval is a click, and nothing ties what the agent does
afterwards to what was approved: it may repeat a step, leave the order, or widen the scope,
and a part of the agent to which work is handed may do all of it. We make a plan a type. A
plan is a regular expression over actions with predicates on their payloads. The approval of
a plan is a value that cannot be copied and whose type is the plan. Performing an action
replaces the plan in the type by its derivative, which is what is left of it. A token for a
plan serves for every plan that it contains, and two parts of a program may share a plan as
long as their steps, interleaved, stay within it. We prove that what happens under an
approval is always the beginning of something that was approved, on every behavior of the
model, across parts of a program, and for code that the model writes and that is checked
before it runs. A token that such code gives back certifies what the code did not do. We
prove that narrowing a plan is sound exactly for containment, so that a change of plan needs
a new question if and only if the new plan is not contained in what is left of the old one,
and the type checker decides which. The single approvals and budgets of an earlier paper are
plans. A script implements the calculus and tests the theorems on \exPrograms{} random
programs, checks the containment theorem on \exPairs{} pairs of plans, and refutes each of
\exRules{} weakened rules by a program.
\end{abstract}

\keywords{regular expressions, derivatives, typestate, linear types, delegation, agents, human oversight}

\authorsaddresses{}

\maketitle

\section{Introduction}

An agent that is asked to make a release does not start typing. It reads, and then it says
what it will do: edit these files, run the tests, push to this branch. A person reads the
plan and approves it. From then on the agent works, and the person looks away. This is how
the harnesses for coding agents work today, and it is the right division of labor: a plan is
short, and a person can judge it.

The approval is worth less than it looks. The plan is prose, and the approval is a click, so
that nothing connects them to what follows.

\emph{Order and number.} The plan says that the tests run before the push. The agent pushes,
finds a problem, edits, and pushes again. Each action occurs in the plan. The sequence does
not.

\emph{Scope.} The plan names a branch and a set of files. The agent, three steps later, is
working from what a tool returned, and the file that it edits is another one.

\emph{Parts.} The agent hands the edits to a helper, or to code that it writes on the spot.
What may the helper do? If it inherits the approval, it may also push.

\emph{Change.} Plans change, because the world answers. When the agent wants to do something
else, it either asks again, and the person is interrupted for a detail, or it does not, and
the approval covers something that was never shown.

An earlier paper of this series treated a single approval as a value that cannot be
copied, with a type that records one action and a predicate on its payload, and a budget as
a value whose type records an amount~\citep{tpp4}. A plan is neither: it is about several
actions and their order. But it has the same shape. This paper generalizes both to one
notion.

\paragraph{A plan is a language.} A \emph{plan} is a regular expression whose letters are
actions with a predicate on the payload:
\[
\edit(x.\,x \le 2)^{*} \cdot \test \cdot \push(x.\,x = 1)
\]
allows any number of edits of the files up to number $2$, then a test, then one push to
branch $1$. The approval of a plan $P$ is a token of type $\Plan{P}$. An action under the
token is well typed if the plan has a word that begins with it, and afterwards the token has
the type of what is left: the derivative of the plan by the action~\citep{brzozowski1964}.
After the push the plan is used up, and a second push does not type. The model may choose
which file to edit, and the type checker asks only that the program has shown the number to
be at most $2$.

\paragraph{Containment.} A token for a plan serves for any plan that it contains. So an
agent may always do less, and may hand a part of the program a smaller plan. A change of
plan needs a new question exactly if the new plan is not contained in what is left of the
old one. Containment of regular languages is decidable, so the type checker settles for
every change of plan in a program whether the person has to be asked.

\paragraph{Parts.} Two parts of a program may each get a token, for plans $P_1$ and $P_2$,
if every interleaving of their steps stays within the plan that is split. For plans that
prescribe an order this is rarely so, and then the token is \emph{lent}: the helper gets the
token and has to give back a token for a stated plan. The type of what comes back says what
the helper has not done.

\paragraph{Results.} As in the earlier paper, a semantics with named tokens and a ledger
serves as the specification, and a semantics without either is what runs.
\begin{itemize}
\item \emph{Safety and erasure} (Theorems~\ref{thm:safety} and~\ref{thm:erasure}). The
  ledger never refuses a step of a well-typed program, so tokens need not exist at run time.
\item \emph{Consent} (Theorem~\ref{thm:consent}). In every execution, the actions under a
  grant form a word that begins a word of the granted plan, whichever parts of the program
  performed them.
\item \emph{Authority and lending} (Theorem~\ref{thm:authority},
  Corollary~\ref{cor:lending}). A part of a program stays within the plans in the types of
  its tokens. If code gives back a token for a plan $Q$, then what it did, followed by
  anything that $Q$ allows, is within the plan that it was lent.
\item \emph{Narrowing is containment} (Theorem~\ref{thm:containment}). The rule is sound,
  and for every pair of plans without containment a program shows that no more could be
  allowed.
\item \emph{Approvals and budgets are plans} (Proposition~\ref{prop:instances}).
\end{itemize}
The artifact tests each statement (Section~\ref{sec:tests}).

\section{One Plan, Several Programs}
\label{sec:overview}

Let $R$ be the plan of the introduction, with $E = \edit(x.\,x \le 2)$:
\[
R \;=\; E^{*} \cdot \test \cdot \push(x.\,x = 1).
\]
A program asks for it and follows it.
\begin{prog}
\letin{r}{\ask{R}}{\kw{case}\;r\;\kw{of}\;\inl y \Rightarrow} & \text{the person approves $R$}\\
\quad \letin{m}{\model}{\kw{if}\;m \le 2\;\kw{then}} & \text{the model names a file}\\
\quad \kw{act}\;\edit\;m\;y\;\kw{as}\;(x_1, y_1)\;\kw{in} & \text{$y_1$ has the plan $R / \edit(m)$}\\
\quad \kw{act}\;\test\;0\;y_1\;\kw{as}\;(x_2, y_2)\;\kw{in} & \text{the plan is now the push}\\
\quad \kw{act}\;\push\;1\;y_2\;\kw{as}\;(x_3, y_3)\;\kw{in}\;\ret 0 & \text{nothing is left}
\end{prog}
The branches for a refusal and for a larger number return at once and are not shown.
The expression $\kw{act}\;a\;t\;y\;\kw{as}\;(x, y')\;\kw{in}\;e$ performs the action $a$ with
the payload $t$ under the token $y$. In $e$, the variable $x$ is the reply of the world, and
$y'$ is the token again, at the type of what is left of its plan. The old name $y$ is gone.
What is left of $R$ after an edit is $R$ itself, after the test it is the push, and after the
push it is the empty plan.

\begin{table}[t]
\caption{Variants of the program of Section~\ref{sec:overview}, code for the type
$\Plan{R} \multimap \Plan{R}$, and what the checker of the artifact says. A verdict other
than ``accepted'' is the reason for which the checker rejects.}
\label{tab:verdicts}
\small
\begin{tabular}{@{}ll@{}}
\toprule
Program & Verdict \\
\midrule
follows the plan & accepted \\
edits the file that the model names, without the test $m \le 2$ & an action that the plan does not allow \\
pushes a second time & an action that the plan does not allow \\
edits once more after the test & an action that the plan does not allow \\
pushes to branch $2$ & an action that the plan does not allow \\
lends the token to code and goes on with what comes back & accepted \\
splits the token into the edits and the rest & parts that go beyond the plan \\
\midrule
code: edits twice and gives the token back & accepted \\
code: runs the test and gives the token back & a result of another type \\
code: edits and keeps the token & a result of another type \\
\bottomrule
\end{tabular}
\end{table}

Table~\ref{tab:verdicts} lists what else one might write. The four rejected programs in its
upper part are the gaps of the introduction. Without the test of $m$, nothing says that the
file is one of the approved ones. A second push, an edit after the test, and a push to
another branch are actions that the plan knows, at a moment or with a payload at which it
does not allow them.

\paragraph{Handing work on.} Suppose the edits are to be made by code that the model writes.
The first idea is to split the token: one part for the edits, $E^{*}$, and one for the rest,
$\test \cdot \push(x.\,x = 1)$. The checker refuses. Two parts of a program act
independently, so every interleaving of their steps has to be within $R$, and the
interleaving that tests first and edits afterwards is not. A plan that prescribes an order
cannot be cut into independent pieces.

The token can be lent. The expression $\run{\Plan{R}}{\Plan{R}}{t}{y}$ asks the model for
code that takes a token for $R$ and \emph{gives back} a token for $R$. Code that edits twice
and returns its token has this type, since what is left of $R$ after two edits is $R$. Code
that runs the test cannot return its token at this type: what is left is the push, and a
token for the push does not serve as a token for $R$. The code has no other way to obtain a
token for $R$ than to ask the person. So when the code returns, the caller holds a token for
$R$ and knows, without having read the code, that under this token the code made only edits
(Corollary~\ref{cor:lending}).

\begin{table}[t]
\caption{An execution of the program that lends its token to code which edits twice. Every
event names the token, and the last column is what the ledger holds for the token after the
event.}
\label{tab:trace}
\small
\begin{tabular}{@{}lll@{}}
\toprule
Event & made by & left of the plan \\
\midrule
$\kw{grant}(0, R)$ & the person & $R$ \\
$\kw{act}(0, \edit, 1)$ & the code & $R$ \\
$\kw{act}(0, \edit, 2)$ & the code & $R$ \\
$\kw{act}(0, \test, 0)$ & the program & $\push(x.\,x = 1)$ \\
$\kw{act}(0, \push, 1)$ & the program & $\epsilon$ \\
\bottomrule
\end{tabular}
\end{table}

Table~\ref{tab:trace} shows an execution. The names of tokens and the ledger are devices of
the semantics that states the guarantee. The program that runs has neither, and its log is
the column of events without the names.

\paragraph{Changing the plan.} Suppose the agent decides, after its edits, that it will not
push. It needs no question: the plan that ends after the test is contained in $R$. Suppose it
wants to push to another branch. That plan is not contained in what it holds, the narrowing
does not type, and the program has to contain a question. In both cases the checker has
decided, before anything ran, whether the person is asked.

\section{The Calculus}
\label{sec:calculus}

\subsection{Plans}

\begin{figure}[t]
\[
\begin{array}{@{}l@{\quad}r@{\;}l@{}}
\text{terms} & t ::= & x \mid n \mid t + t \\
\text{predicates} & \varphi ::= & \top \mid t \le t \mid t = t \mid \neg\varphi \mid \varphi \wedge \varphi \\
\text{plans} & P, Q ::= & 0 \mid \epsilon \mid a(x.\varphi) \mid P \cdot P \mid P + P \mid P^{*} \mid P \parallel P \mid P / a(t) \\
\text{types} & \tau ::= & \Base \mid 1 \mid \tau \otimes \tau \mid \tau \oplus \tau \mid \Plan{P} \\
\text{operands} & w ::= & y \mid t \mid () \mid (w, w) \mid \inl w \mid \inr w \\
\text{expressions} & e ::= & \ret w \mid \letin{x}{e}{e} \mid \letin{y}{e}{e} \mid \letin{(y, y)}{w}{e} \\
 & \mid & \kw{case}\;w\;\kw{of}\;\inl y \Rightarrow e;\ \inr y \Rightarrow e \mid \ite{\varphi}{e}{e} \mid \model \\
 & \mid & \ask{P} \mid \act{a}{t}{w}{x, y}{e} \mid \splt{w}{P}{P} \\
 & \mid & \run{\tau}{\tau}{t}{w} \mid \iter{n}{w}{y}{e}
\end{array}
\]
\Description{The grammar of terms, predicates, plans, types, operands, and expressions.}
\caption{Syntax. The letter $a$ ranges over actions, $n$ over numbers, $x$ over base
variables, and $y$ over resource variables.}
\label{fig:syntax}
\end{figure}

Figure~\ref{fig:syntax} gives the syntax. A \emph{letter} is a pair $a(n)$ of an action and
a number, its payload. A \emph{word} is a sequence of letters. A closed plan denotes a set
$\lang{P}$ of words in the usual way: $0$ is the empty set, $\epsilon$ the set of the empty
word, $a(x.\varphi)$ the set of the letters $a(n)$ for which $\varphi$ holds of $n$;
$P \cdot Q$, $P + Q$ and $P^{*}$ are concatenation, union and iteration; and $P \parallel Q$
is the shuffle, the set of all interleavings of a word of $P$ with a word of $Q$. The plan
$P / a(t)$ is what is left of $P$ after the letter: $\lang{P / a(n)}$ is the set of the
words $w$ with $a(n)\,w \in \lang{P}$. It can be written without the stroke, by the rules of
\citet{brzozowski1964}, which extend to the shuffle by
$(P \parallel Q) / \ell = (P / \ell \parallel Q) + (P \parallel Q / \ell)$. We keep it in
the syntax because types need it for a payload that is not yet known.

A token may be dropped at any time, so what counts is not the words of a plan but their
beginnings.

\begin{definition}[Authority]
\label{def:authority}
The \emph{authority} $\auth{P}$ of a closed plan is the set of the prefixes of the words of
$\lang{P}$, together with the empty word. For sets of words $A$ and $B$, $A \sh B$ is the
set of the interleavings of a word of $A$ with a word of $B$.
\end{definition}

So $\auth{0} = \auth{\epsilon} = \{\epsilon\}$: a token for either plan allows nothing. A
letter $\ell$ is in $\auth{P}$ exactly if $\lang{P / \ell}$ is not empty.

\subsection{Expressions}

Expressions are those of the earlier calculus~\citep{tpp4}, with other operations on tokens.
$\ask{P}$ shows the person the plan $P$, with the values of its variables filled in, and
returns a token or a refusal. $\act{a}{t}{w}{x, y}{e}$ performs an action under a token.
$\splt{w}{P_1}{P_2}$ divides a token in two. $\run{\tau_1}{\tau_2}{t}{w}$ asks the model for
code from $\tau_1$ to $\tau_2$, checks it, and runs it on $w$, and $\iter{n}{w}{y}{e}$ is a
bounded loop with the state $y$. The model's answers and the world's replies are numbers.

\subsection{Two Semantics}
\label{sec:semantics}

\begin{figure}[t]
\begin{mathpar}
\inferrule*[right=G]{k \notin \dom L}{L \lstep{\kw{grant}(k, P)} L[k \mapsto P]}
\and
\inferrule*[right=A]{a(n) \in \auth{L(k)}}{L \lstep{\kw{act}(k, a, n)} L[k \mapsto L(k) / a(n)]}
\and
\inferrule*[right=S]{\auth{P_1} \sh \auth{P_2} \subseteq \auth{L(k)} \\ k' \notin \dom L}
  {L \lstep{\kw{split}(k, k', P_1, P_2)} L[k' \mapsto P_1,\, k \mapsto P_2]}
\end{mathpar}
\vspace{-1ex}
\begin{mathpar}
\inferrule*[right=Ask-Yes]{L \lstep{\kw{grant}(k, P)} L'}{L, \ask{P} \evto L', \inl\tok{k}, \kw{grant}(k, P)}
\and
\inferrule*[right=Ask-No]{ }{L, \ask{P} \evto L, \inr(), \kw{deny}}
\and
\inferrule*[right=Act]{\sem{w} = \tok{k} \\ L \lstep{\kw{act}(k, a, \sem{t})} L_1 \\ L_1, e[n/x, \tok{k}/y] \evto L_2, v, s}
  {L, \act{a}{t}{w}{x, y}{e} \evto L_2, v, \kw{act}(k, a, \sem{t}) \cdot s}
\and
\inferrule*[right=Split]{\sem{w} = \tok{k} \\ L \lstep{\kw{split}(k, k', P_1, P_2)} L'}
  {L, \splt{w}{P_1}{P_2} \evto L', (\tok{k'}, \tok{k}), \kw{split}(k, k', P_1, P_2)}
\end{mathpar}
\Description{Three rules for the steps of a ledger and four rules of evaluation.}
\caption{The steps of a ledger, and the evaluation of the three operations on tokens. In
\textsc{Act}, $n$ is any number: the reply of the world. The other rules are those of
\citet{tpp4}.}
\label{fig:semantics}
\end{figure}

In the first semantics a token has a name $k$, and a \emph{ledger} $L$ maps each name to a
closed plan: what is left for this token. Figure~\ref{fig:semantics} shows its steps. A
grant enters a plan. An action is possible if its letter is in the authority of the plan,
and it leaves the derivative. A split replaces a plan by two whose interleavings it
contains. The judgment $L, e \evto L', v, s$ evaluates a closed expression to a value and a
trace of events; it goes \emph{wrong} if the ledger has no step for an action or a split, or
if a value has the wrong shape. The rules for the other expressions are as in the earlier
paper. In particular $\kw{run}$ runs the code that the model returns if a check
$y{:}\tau_1 \vdash_{\kw{c}} e : \tau_2$ accepts it, with the result in the left summand, and
otherwise returns its argument in the right one.

The second semantics, $e \evto^{\circ} v, s$, has no ledger. A token is a value $\bullet$
without content, a split yields two of them, and the events lose their names. It performs
whatever action the program asks for.

\section{Types}
\label{sec:types}

\begin{figure}[t]
\begin{mathpar}
\inferrule*[right=Ask]{ }{\Delta \vdash \ask{P} : \Plan{P} \oplus 1}
\and
\inferrule*[right=Narrow]{\Delta \vdash w : \Plan{P} \\ \Phi \models \auth{Q} \subseteq \auth{P}}
  {\Delta \vdash w : \Plan{Q}}
\and
\inferrule*[right=Act]{\Delta_1 \vdash w : \Plan{P} \\ \Phi \models a(t) \in \auth{P} \\
  \Phi; \Gamma, x; \Delta_2, y{:}\Plan{P / a(t)} \vdash e : \tau \\ x \notin \Gamma \cup \mathrm{fv}(\tau)}
  {\Delta_1, \Delta_2 \vdash \act{a}{t}{w}{x, y}{e} : \tau}
\and
\inferrule*[right=Split]{\Delta \vdash w : \Plan{P} \\ \Phi \models \auth{P_1} \sh \auth{P_2} \subseteq \auth{P}}
  {\Delta \vdash \splt{w}{P_1}{P_2} : \Plan{P_1} \otimes \Plan{P_2}}
\and
\inferrule*[right=Run]{\Delta \vdash w : \tau_1 \\ \tau_1, \tau_2\ \text{closed}}
  {\Delta \vdash \run{\tau_1}{\tau_2}{t}{w} : \tau_2 \oplus \tau_1}
\and
\inferrule*[right=Iter]{\Delta \vdash w : \tau \\ \Phi; \Gamma; y{:}\tau \vdash e : \tau}
  {\Delta \vdash \iter{n}{w}{y}{e} : \tau}
\end{mathpar}
\Description{Six typing rules for the operations on tokens, code and loops.}
\caption{Typing of the operations on tokens. A judgment $\Delta \vdash e : \tau$ stands for
$\Phi; \Gamma; \Delta \vdash e : \tau$. The premises with $\Phi \models$ hold if the stated
relation holds for the closed plans under every assignment of numbers to the variables that
satisfies $\Phi$. The other rules are those of \citet{tpp4}: contexts are divided among the
parts of an expression, the branches of a conditional get its test as a fact, and
\textsc{Narrow} extends to pairs and sums.}
\label{fig:typing}
\end{figure}

\begin{figure}[t]
\begin{mathpar}
\inferrule*[right=Var]{ }{\Delta, y{:}\tau \vdash y : \tau}
\and
\inferrule*[right=Tm]{ }{\Delta \vdash t : \Base}
\and
\inferrule*[right=Pair]{\Delta_1 \vdash w_1 : \tau_1 \\ \Delta_2 \vdash w_2 : \tau_2}
  {\Delta_1, \Delta_2 \vdash (w_1, w_2) : \tau_1 \otimes \tau_2}
\and
\inferrule*[right=Inl]{\Delta \vdash w : \tau_1}{\Delta \vdash \inl w : \tau_1 \oplus \tau_2}
\and
\inferrule*[right=Model]{ }{\Delta \vdash \model : \Base}
\and
\inferrule*[right=Let-B]{\Delta_1 \vdash e_1 : \Base \\ \Phi; \Gamma, x; \Delta_2 \vdash e_2 : \tau \\
  x \notin \Gamma \cup \mathrm{fv}(\tau)}
  {\Delta_1, \Delta_2 \vdash \letin{x}{e_1}{e_2} : \tau}
\and
\inferrule*[right=Let-R]{\Delta_1 \vdash e_1 : \tau_1 \\ \Delta_2, y{:}\tau_1 \vdash e_2 : \tau}
  {\Delta_1, \Delta_2 \vdash \letin{y}{e_1}{e_2} : \tau}
\and
\inferrule*[right=Let-P]{\Delta_1 \vdash w : \tau_1 \otimes \tau_2 \\ \Delta_2, y_1{:}\tau_1, y_2{:}\tau_2 \vdash e : \tau}
  {\Delta_1, \Delta_2 \vdash \letin{(y_1, y_2)}{w}{e} : \tau}
\and
\inferrule*[right=Case]{\Delta_1 \vdash w : \tau_1 \oplus \tau_2 \\ \Delta_2, y_1{:}\tau_1 \vdash e_1 : \tau \\
  \Delta_2, y_2{:}\tau_2 \vdash e_2 : \tau}
  {\Delta_1, \Delta_2 \vdash \kw{case}\;w\;\kw{of}\;\inl y_1 \Rightarrow e_1;\ \inr y_2 \Rightarrow e_2 : \tau}
\and
\inferrule*[right=If]{\Phi, \varphi; \Gamma; \Delta \vdash e_1 : \tau \\ \Phi, \neg\varphi; \Gamma; \Delta \vdash e_2 : \tau}
  {\Delta \vdash \ite{\varphi}{e_1}{e_2} : \tau}
\end{mathpar}
\Description{Ten typing rules for operands, bindings, sums and conditionals.}
\caption{The rules that do not concern tokens, from \citet{tpp4}. A context is divided among
the parts of an expression, and the two branches of a case or a conditional share one. There
are the expected rules for $()$, $\inr$ and $\kw{ret}$.}
\label{fig:rest}
\end{figure}

The judgment $\Phi; \Gamma; \Delta \vdash e : \tau$ has the three contexts of the earlier
paper: facts, base variables, and resource variables with their types.
Figure~\ref{fig:rest} repeats the rules that have nothing to do with plans. A rule with
several parts divides $\Delta$ among them, so a resource variable is used at most once, and
it may be left unused. A conditional adds its test to the facts of one branch and the
negation to those of the other; this is how a program shows that a payload satisfies the
predicate of a letter. Figure~\ref{fig:typing} shows the rules that are new.

\paragraph{Acting.} \textsc{Act} needs a token for a plan in whose authority the action
lies, whatever the values of the variables, as far as the facts say. The continuation gets
the token at the type $\Plan{P / a(t)}$. The stroke stays in the type until the payload is
known; in the example of Section~\ref{sec:overview} the facts $m \le 2$ suffice to see that
$R / \edit(m)$ and $R$ have the same authority.

\paragraph{Narrowing.} \textsc{Narrow} is subtyping for tokens: a token for $P$ serves as a
token for $Q$ if the authority of $Q$ is contained in that of $P$.

\paragraph{Splitting.} \textsc{Split} asks that every interleaving of what the two parts may
do lies in what the whole may do. A plan $P_1 \parallel P_2$ splits into $P_1$ and $P_2$, a
plan $P^{*}$ into two copies of itself, and every plan into itself and $\epsilon$.

\paragraph{Loops.} The state of a loop has one type for all turns. For a token this means a
plan that still holds after what a turn does, and such plans are the starred ones: a turn
that performs a word of $P$ under a token for $P^{*}$ leaves a plan that contains $P^{*}$,
and \textsc{Narrow} restores the type. The star of a plan is the type of a loop.

\paragraph{Deciding the premises.} For closed plans, membership of a letter is the test
whether a derivative denotes the empty set, and the two containments are instances of the
inclusion problem for regular expressions with shuffle. Both are decidable, the second at a
cost that grows fast with the size of the plans~\citep{stockmeyer1973,mayer1994}. Plans that
a person is to read are small. With variables, the premises quantify over the assignments
that satisfy the facts; the predicates are those of linear arithmetic, and the letters of a
plan can be divided into finitely many classes on which all its predicates agree, as for
symbolic automata~\citep{dantoni2017}. The artifact works with four numbers and three
actions and tries all assignments.

\section{Results}
\label{sec:results}

\subsection{Derivatives and Authority}

\begin{lemma}[Derivatives]
\label{lem:derivatives}
Let $P$, $Q$, $P_1$, $P_2$ be closed plans and $\ell$ a letter.
\begin{enumerate}
\item $\ell\,w \in \auth{P}$ if and only if $\ell \in \auth{P}$ and $w \in \auth{P / \ell}$.
\item If $\auth{Q} \subseteq \auth{P}$ and $\ell \in \auth{Q}$, then
  $\auth{Q / \ell} \subseteq \auth{P / \ell}$.
\item If $\auth{P_1} \sh \auth{P_2} \subseteq \auth{P}$ and $\ell \in \auth{P_1}$, then
  $\auth{P_1 / \ell} \sh \auth{P_2} \subseteq \auth{P / \ell}$.
\end{enumerate}
\end{lemma}

\begin{proof}
(1) A word $\ell\,w$ is a prefix of a word of $\lang{P}$ if and only if $w$ is a prefix of a
word of $\lang{P / \ell}$; the empty word is in $\auth{P / \ell}$ in any case, which is why
the condition $\ell \in \auth{P}$ is stated. (2) Let $w \in \auth{Q / \ell}$. By (1),
$\ell\,w \in \auth{Q} \subseteq \auth{P}$, and by (1) again $w \in \auth{P / \ell}$. (3) An
interleaving of $w_1 \in \auth{P_1 / \ell}$ and $w_2 \in \auth{P_2}$, with $\ell$ in front,
is an interleaving of $\ell\,w_1 \in \auth{P_1}$ and $w_2$, hence in $\auth{P}$, and (1)
applies.
\end{proof}

\subsection{Safety and Erasure}

A value $\tok{k}$ has the closed type $\Plan{P}$ in a ledger $L$, and owns the name $k$, if
$\auth{P} \subseteq \auth{L(k)}$: the ledger allows at least what the type promises. Pairs
own disjoint sets of names, and the judgment $L \vdash v : \tau \own O$ and its extension to
environments $\gamma$ are as in the earlier paper.

\begin{theorem}[Safety]
\label{thm:safety}
Let $\Phi; \Gamma; \Delta \vdash e : \tau$ and $L \vdash \gamma : \Phi; \Gamma; \Delta \own O$.
Then no evaluation of $\gamma(e)$ from $L$ goes wrong. If $L, \gamma(e) \evto L', v, s$ and
$N$ is the set of the names of $L'$ that are not in $L$, then
$L' \vdash v : \gamma(\tau) \own O'$ for some $O' \subseteq O \cup N$; $L'$ agrees with $L$
on the names of $L$ outside $O$; and every name in an event of $s$ is in $O \cup N$.
\end{theorem}

\begin{proof}
By induction on the derivation, first without \textsc{Run} and then with it, as for the
corresponding theorem of~\citet{tpp4}. The cases that are new:

\emph{Act.} The operand is a token $\tok{k}$ with $k \in O$ and
$\auth{\gamma(P)} \subseteq \auth{L(k)}$. The numbers of $\gamma$ satisfy $\Phi$, so the
letter $\ell = a(\sem{\gamma(t)})$ is in $\auth{\gamma(P)}$ and hence in $\auth{L(k)}$: the
ledger has the step. Afterwards the entry is $L(k) / \ell$, and by
Lemma~\ref{lem:derivatives}(2) the token has the type $\Plan{\gamma(P) / \ell}$, which is
the type of $y$ under the environment extended by the reply. No other entry has changed, so
the rest of the environment keeps its types, and the induction hypothesis applies to the
continuation.

\emph{Narrow.} If $\auth{\gamma(Q)} \subseteq \auth{\gamma(P)} \subseteq \auth{L(k)}$, the
token has the type $\Plan{\gamma(Q)}$.

\emph{Split.} From the premise and the type of the token,
$\auth{\gamma(P_1)} \sh \auth{\gamma(P_2)} \subseteq \auth{L(k)}$, so the ledger has the
step. It enters $\gamma(P_1)$ under a new name and $\gamma(P_2)$ under $k$, and the two
tokens have exactly the types of the result.
\end{proof}

\begin{theorem}[Erasure]
\label{thm:erasure}
Let $\cdot; \cdot; \cdot \vdash e : \tau$. Then $e \evto^{\circ} u, r$ if and only if
$\emptyset, e \evto L, v, s$ for some $L$, $v$, and $s$ that are $u$ and $r$ once the names
and the split events are removed.
\end{theorem}

\begin{proof}
No rule consults the ledger except to ask whether a step exists, and no value depends on a
name. So the two semantics have the same evaluations up to the ones that go wrong in the
first, and by Theorem~\ref{thm:safety} there are none.
\end{proof}

\subsection{What Happens Was Approved}

\begin{definition}[Explained traces]
\label{def:explained}
In a trace, a name is \emph{under} a grant $\kw{grant}(g, P)$ if it is $g$ or is introduced
by a split of a name under the grant. The trace is \emph{explained} if every name is
introduced once, before any other event that contains it, and for every grant
$\kw{grant}(g, P)$ the actions on the names under it, in the order of the trace, form a word
of $\auth{P}$.
\end{definition}

The definition does not mention how a plan was divided. Whatever the parts did, taken
together it has to be the beginning of a word of the plan that the person saw.

\begin{lemma}[The ledger explains]
\label{lem:ledger}
If $\emptyset \lstep{s} L$, then $s$ is explained.
\end{lemma}

\begin{proof}
For a grant of $P$ under $g$, let $u$ be the word of the actions so far on names under it,
and let $k_1, \ldots, k_m$ be these names. We show that $u \in \auth{P}$ and
\[
\auth{L(k_1)} \sh \cdots \sh \auth{L(k_m)} \;\subseteq\; \auth{P / u}.
\]
After the grant, $m = 1$ and both sides are $\auth{P}$. A split replaces one $\auth{L(k_i)}$
by a set $\auth{P_1} \sh \auth{P_2}$ that it contains, and the shuffle is monotone. An
action $\ell$ on $k_i$ needs $\ell \in \auth{L(k_i)}$. Each of the sets contains the empty
word, so $\ell$ is in the left side and hence in $\auth{P / u}$, which gives
$u\,\ell \in \auth{P}$ by Lemma~\ref{lem:derivatives}(1). The inclusion for $u\,\ell$ is
Lemma~\ref{lem:derivatives}(3), for $m$ parts.
\end{proof}

\begin{theorem}[Consent]
\label{thm:consent}
Let $\cdot; \cdot; \cdot \vdash e : \tau$. If $e \evto^{\circ} u, r$, then $r$ is an
explained trace without its names and split events.
\end{theorem}

\begin{proof}
By Theorem~\ref{thm:erasure} and Lemma~\ref{lem:ledger}, since every evaluation with a ledger
takes the steps of its trace.
\end{proof}

In words: in every execution of a well-typed program, without tokens and without a ledger,
the actions can be assigned to the grants of the person such that the actions of each grant,
in the order in which they happened, are the beginning of a word of the plan that was shown.
This holds for every answer of the model and every reply of the world.

\subsection{A Part Stays within Its Types}

For a value $v$ of a closed type $\tau$, the \emph{least ledger} $\least{v}{\tau}$ maps the
name of every token in $v$ to the plan in its type.

\begin{theorem}[Authority]
\label{thm:authority}
Let $\Phi; \Gamma; \Delta \vdash e : \tau$ and $L \vdash \gamma : \Phi; \Gamma; \Delta \own O$,
and let $L, \gamma(e) \evto L', v, s$. Then the least ledger of $\gamma$ at the types of
$\Delta$ has the steps of $s$.
\end{theorem}

\begin{proof}
As in the earlier paper: the least ledger gives $\gamma$ its types, an evaluation depends on
the ledger only through the existence of steps, and by Theorem~\ref{thm:safety} the
evaluation from the least ledger does not go wrong.
\end{proof}

With Lemma~\ref{lem:ledger}, applied to the least ledger as a sequence of grants, the theorem
says that the actions of a part under a token of type $\Plan{P}$, and under everything that
is split from it, form a word of $\auth{P}$: the plan of the type, also where the ledger
holds more. Every other action of the part is under a grant that the person made while the
part ran.

\begin{corollary}[Lending]
\label{cor:lending}
Let code be accepted for $\Plan{P} \multimap \Plan{Q}$ and run on a token $\tok{k}$, and let
$u$ be the word of its actions on the names under $k$. If the code returns a token whose
name is under $k$, then $u\,w \in \auth{P}$ for every $w \in \auth{Q}$.
\end{corollary}

\begin{proof}
Evaluate from the least ledger $[k \mapsto P]$. By Theorem~\ref{thm:safety} the result
$\tok{k'}$ has the type $\Plan{Q}$ in the final ledger $L'$, so
$\auth{Q} \subseteq \auth{L'(k')}$. By the invariant in the proof of Lemma~\ref{lem:ledger},
$\auth{L'(k')}$ is contained in the shuffle of the entries under $k$, which is contained in
$\auth{P / u}$. So $w \in \auth{P / u}$, and $u \in \auth{P}$.
\end{proof}

The corollary is the reason why lending works for plans with an order. The borrower cannot
make a token, so a token for $Q$ that comes back is the one that was lent, or a part of it,
or one for which the person was asked again. In the first two cases its type is a statement
about the past: nothing was done that would make $Q$ impossible. In the example of
Section~\ref{sec:overview}, with $Q = P = R$, no test and no push can have happened.

\subsection{Narrowing Is Containment}

\begin{theorem}[Containment]
\label{thm:containment}
Let $P$ and $Q$ be closed plans.
\begin{enumerate}
\item If $\auth{Q} \subseteq \auth{P}$, a token for $P$ may be used as a token for $Q$, and
  the theorems of this section hold.
\item If $\auth{Q} \not\subseteq \auth{P}$, there is a program that is well typed if a token
  for $P$ may be used as a token for $Q$, and whose evaluation from the empty ledger goes
  wrong.
\end{enumerate}
\end{theorem}

\begin{proof}
(1) is the rule \textsc{Narrow}. (2) Take a word $u \in \auth{Q}$ that is not in $\auth{P}$.
The program asks for $P$, uses the token as a token for $Q$, and performs the letters of $u$
one after the other. Each action is well typed by Lemma~\ref{lem:derivatives}(1), since every
prefix of $u$ is in $\auth{Q}$. If the person grants, the ledger holds $P$, and it has no
step for the first letter with which $u$ leaves $\auth{P}$.
\end{proof}

So the question of the introduction has an exact answer. An agent that holds a token for $P$
and wants to proceed by a plan $Q$ has to ask if and only if $\auth{Q}$ is not contained in
$\auth{P}$. If it is contained, asking is unnecessary. If it is not, then by (2) no sound
rule could let the program go on without a question. The type checker decides, since a
program that narrows where it should ask does not type.

\subsection{Approvals and Budgets Are Plans}

The earlier paper has two kinds of tokens~\citep{tpp4}. An approval for an action $a$ with a
payload that satisfies $\varphi$ is used up by one action. A budget for $a$ of the amount $n$
is used up by payloads that sum to $n$, and it splits into budgets whose amounts add up to
it.

\begin{proposition}[Instances]
\label{prop:instances}
Let payloads be the numbers up to some bound. An approval is a token for the plan
$a(x.\varphi)$. A budget of $n$ is a token for the plan $B_n$, where
\[
B_n \;=\; a(x.\,x = 0)^{*} \cdot \Bigl(\epsilon + \sum_{m = 1}^{n} a(x.\,x = m) \cdot B_{n - m}\Bigr).
\]
The authority of $B_n$ is the set of the words of letters $a(m)$ whose payloads sum to at
most $n$. For $m \le n$, the plan $B_n / a(m)$ has the authority of $B_{n - m}$, and for
$m > n$ the letter $a(m)$ is not in $\auth{B_n}$. And
$\auth{B_s} \sh \auth{B_t} \subseteq \auth{B_n}$ if and only if $s + t \le n$.
\end{proposition}

\begin{proof}
The description of the authority is by induction on $n$. The other claims follow from it: a
word whose payloads sum to at most $n - m$ can follow $a(m)$, and an interleaving of two
words has the sum of their sums.
\end{proof}

So the rules of the earlier paper for spending and splitting a budget are \textsc{Act} and
\textsc{Split} for $B_n$, and its rule for narrowing an approval is \textsc{Narrow} for a
letter. What plans add is everything between one action and an amount: order, alternatives,
repetition of a group of actions, and bounds on some actions and not on others.

\section{Plans in Use}
\label{sec:use}

\subsection{Three Plans}

The operators of regular expressions correspond to things that people say when they approve
a plan. Three examples, each checked by the artifact with a program that keeps the plan and
one that does not.

\paragraph{A condition on the order.} ``Push as often as you like, but never without a test
since the last edit'':
\[
(\edit^{*} \cdot \test \cdot \push)^{*}.
\]
A program that edits, tests, pushes, edits twice, tests, and pushes is accepted. One that
edits and pushes right away is rejected at the push: what is left of the plan after the edit
begins with an edit or a test.

\paragraph{A bound on one action.} ``Edit and test freely; push at most twice'':
\[
(\edit + \test)^{*} \;\parallel\; (\push + \epsilon) \cdot (\push + \epsilon).
\]
The shuffle says that the two pushes may come anywhere among the other actions. Without it
the plan would need a copy of the starred part before, between, and after the pushes. A
third push is rejected. This is a budget for one action inside a plan that leaves the others
free.

\paragraph{Two helpers.} ``One of you takes the files up to number 1, the other takes file
2 and runs the tests'':
\[
\edit(x.\,x \le 1)^{*} \;\parallel\; (\edit(x.\,x = 2) + \test)^{*}.
\]
A token for this plan splits into two, one for each side, by the rule \textsc{Split}, and
the two parts of the program proceed independently. If the first helper edits file 2, its
action is rejected against its own half. A split of the first plan above into its edits and
the rest is rejected, as in Section~\ref{sec:overview}: what a plan with an order allows to
two independent parts is less than what it allows to one.

\subsection{Asking Again}

By Theorem~\ref{thm:containment}, the places of a program at which the plan changes are of
two kinds. Where the new plan is contained in what the program holds, a narrowing types, and
nothing happens at run time: no event, no question. Where it is not, the program does not
type unless it asks. A person who approves a plan is therefore asked again only when the
agent wants more than was shown, and is never asked when it wants less.

The \exPrograms{} programs of the tests contain \exNarrowings{} narrowings and
\exQuestions{} questions. These are random programs, and the proportion says nothing about
real agents. It does show the two kinds side by side in programs that the checker accepts,
and the test of Theorem~\ref{thm:containment} shows that none of the narrowings could have
been a question that was skipped.

How large a plan should be is a trade. A larger plan means fewer questions and more that can
go wrong without one: the authority of a plan is exactly what its holder may do without the
person hearing of it. A question has a price as well, in the time of the person and in what
waiting does to the cost of the agent~\citep{tpp6}. The calculus does not choose. It makes
the choice an explicit object, the plan in the grant, and it guarantees that the object means
what it says.

\section{Testing the Claims}
\label{sec:tests}

The artifact is one Python script without dependencies. It contains regular expressions with
shuffle and their derivatives, the search that decides the two containments, the type
checker, both semantics, Definition~\ref{def:explained} as a function that does not use the
ledger, a generator of plans and programs, and the tests. All numbers are those of one fixed
seed, and \texttt{make -C artifact check} computes them again and compares.

\paragraph{The expressions.} Everything rests on derivatives, so the script first compares
them with the definition of a language. For \exEnginePlans{} random plans it checks, on
\exEngineWords{} words in all, that a word is in the language by derivatives exactly if it
is by the definition, with the shuffle computed by trying all ways to divide the word.

\paragraph{The theorems.} The generator builds programs that follow the typing rules: it
keeps the contexts, asks for random plans, and searches for actions that the plan at hand
allows under the facts. It narrows, splits, loops under starred plans, hands tokens to code,
and lends them. With a switch it \emph{slips}: at each choice that a rule constrains it
ignores the rule with probability \exSlip{}. We took \exTyped{} programs without slips and
\exTyped{} with slips that the checker accepts, of which \exLucky{} contain a slip that did
no harm. Each program ran three times with the ledger, and again without tokens on the same
random stream. In every execution the ledger refuses no step
(Theorem~\ref{thm:safety}); the execution without tokens gives the same value and trace
(Theorem~\ref{thm:erasure}); the trace is explained (Theorem~\ref{thm:consent}); and wherever
accepted code runs and at every turn of a loop, the part changes no entry that it does not
own and the least ledger of its argument has the steps of its trace
(Theorem~\ref{thm:authority}). The \exRuns{} executions contain \exAsks{} questions,
\exGrants{} grants, \exActs{} actions and \exSplits{} splits. The model wrote \exWritten{}
pieces of code, half of them for the type that was asked and the others for a larger plan
or with slips; \exAccepted{} were accepted. In \exLent{} runs of accepted code a token came
back that was under the one that was lent, and in each the statement of
Corollary~\ref{cor:lending} holds: what the returned type allows is within what is left of
the lent plan after what the code did.

\paragraph{Containment.} For Theorem~\ref{thm:containment} the script draws \exPairs{} pairs
of closed plans. In \exPairsYes{} of them the first contains the second; it then takes a
random walk in the authority of the second and checks that the program which performs it
under a narrowed token is well typed and is not refused. In the other \exPairsNo{} the search
returns a word. The script checks that the word is in the authority of the second plan and
not of the first, that the program of the proof is rejected by the checker, and that it is
accepted and then refused by the ledger when the rule \textsc{Narrow} is replaced by one that
always applies.

\paragraph{Instances.} For Proposition~\ref{prop:instances} the script builds the plans
$B_n$ over its four numbers and checks the three claims for every amount, every payload, and
every way to divide an amount in two.

\begin{table}[t]
\caption{Ten checkers with one weakened rule. Each accepts a program that the ledger
refuses in one of three executions. The last column is the number of random programs that
were tried until the first such program.}
\label{tab:broken}
\small
\begin{tabular}{@{}llr@{}}
\toprule
Rule & What the weakened checker allows & Tried \\
\midrule
\textsc{Act} & an action that the plan does not allow & \exTryAct \\
\textsc{Act} & the token keeps its plan after an action & \exTryResidual \\
\textsc{Narrow} & a token serves for any plan & \exTryNarrow \\
\textsc{Split} & any two parts & \exTrySplit \\
\textsc{Run} & code runs without the check & \exTryRun \\
division of $\Delta$ & a token is still there after it was used & \exTryContraction \\
conditionals & a token counts as used only if both branches use it & \exTryBranch \\
\textsc{Iter} & the body of a loop uses a token from outside & \exTryLoop \\
conditionals & the second branch assumes the test and not its negation & \exTryIfFact \\
bindings & a base variable is bound a second time & \exTryShadow \\
\bottomrule
\end{tabular}
\end{table}

\paragraph{Every rule is needed.} Table~\ref{tab:broken} weakens the checker in \exRules{}
ways. For each the script draws programs from the generator with slips, here with
probability \exSlipBroken{}, until the weakened checker accepts one that the ledger refuses.
The second row is the rule that is specific to plans: if an action left the plan of its token
as it was, a plan with one push would allow any number. The slowest to refute is the rule for
branches, after \exSlowest{} programs.

\paragraph{What the tests do not show.} The plans are small, the payloads are four numbers,
and the programs have a nesting depth of at most six. The containment search is exact on
such plans; it has a bound on its size, beyond which it refuses, and the bound was not
reached in the tests of Theorem~\ref{thm:containment}, where a refusal without a word would
have failed a check.

\section{Limits}
\label{sec:limits}

\paragraph{Replies.} A plan is about what the agent does. It cannot say ``push only if the
tests pass'', because the outcome of the tests is a reply of the world and replies are not
letters. A program can decide by a conditional, and the plan then allows the push in either
case. Putting replies into letters is possible and changes the theory at one point: a part of
a split plan can then become impossible through a reply that the plan did not foresee, and
the shuffle of the parts no longer describes what the others may still do. We have not found
the right definition.

\paragraph{Regular plans.} A plan cannot count without a bound or match an opening with a
closing action at any depth. Budgets are regular only because payloads are bounded.

\paragraph{Permission, not duty.} A token may be dropped, so the theorems speak of
beginnings of words. They do not say that a plan is carried out. A token that must be used up
would be a promise of the agent and not a permission of the person, and it would need a way
for the agent to be released from it.

\paragraph{What the person reads.} A regular expression is a formal object, and a person
approves it only if it is shown in a form that the person can read. The step from the
expression to that form, and back from a plan in prose to an expression, is not part of the
calculus. The theorems are about the expression that is in the grant.

\paragraph{Cost of deciding.} Containment with shuffle is expensive in the worst
case~\citep{mayer1994}. The checker of the artifact enumerates assignments and searches
pairs of derivatives, which is adequate for the plans of the tests and not a decision
procedure that one would ship.

\paragraph{Proofs.} The proofs are on paper, and the artifact tests them and does not prove
them.

\section{Related Work}
\label{sec:related}

\paragraph{Types that follow a protocol.} Typestate lets the operations that a value permits
depend on its state~\citep{strom1986}, and systems such as Vault track the state linearly to
enforce protocols on resources~\citep{deline2001}. \citet{igarashi2002} give types to
resources that carry a \emph{usage}, an expression that describes in which order the
resource will be accessed, and check programs against it; our plans are usages in this
sense, with predicates on payloads, read as permissions that a person grants, and with a
split by shuffle. Session types describe the order of the messages on a
channel~\citep{honda1993}. History effects go the other way: a type and effect system
computes a description of the events that a program may cause, which is then checked against
a policy~\citep{skalka2004}. A plan could be inferred in that way and shown to the person.
The token is needed where the program is not fixed when the plan is approved, which is the
case for code that the model writes.

\paragraph{Derivatives.} What is left of a regular expression after a letter is due to
\citet{brzozowski1964}; \citet{owens2009} show that derivatives make practical automata.
We use them as the transition of a type. Predicates on letters are the idea of symbolic
automata~\citep{dantoni2017}. The inclusion problem for regular expressions, and its cost
with interleaving, are classical~\citep{stockmeyer1973,mayer1994}.

\paragraph{Monitors.} A security automaton accepts the executions that a policy
allows~\citep{schneider2000}. The ledger is one, with a state per token. As in the earlier
paper, the type system shows that the monitor never fires, and the state of the monitor is
held by the program in pieces, which is what makes it possible to give a part of it away.

\paragraph{This series.} The calculus extends that of consent as a linear
resource~\citep{tpp4}, whose two kinds of tokens are instances
(Proposition~\ref{prop:instances}). Which actions need a question at all, given where their
payloads come from, is the subject of~\citet{tpp7}, and what a question costs
of~\citet{tpp6}. This paper adds when a question that was asked still covers what the agent
is about to do. Its artifact follows the practice of~\citet{tpp10}: the quantities of the
text, parameters included, are generated by the script.

\section{Conclusion}

A plan that a person approves is a set of sequences of actions, and the natural name for
such a set is a language. Treating the approval as a value whose type is the language, and
letting the type follow the actions by derivatives, makes three things checkable before an
agent runs: that it stays within the plan, that a part of it stays within a part of the
plan, and that it asks again exactly when the plan grows.

\bibliographystyle{ACM-Reference-Format}
\bibliography{refs}

\end{document}
